
Possible Paper Categories:

4000 words (now “medium”, formerly “short”)

2000 words (now “short”, formerly “demo”)

Structure of the paper:

\section{Introduction}
present the topic and outline the structure of the paper

Including what is FAIR \citet{wilkinson_fair_2016} Open Science \citet{burgelman_open_2019} and how open is not enough \citet{chen_open_2019}.

\section{FAIR, Open Source Digital Musical Interfaces (DMI) at NIME}

Section explores the work that has been done and assessment of 

\subsection{Relevant Work}
Previous papers that have addressed open source hardware and the open sourcing at NIME

Of note:

\begin{enumerate}
    \item  \citet{fiordelmondo_towards_2025}
    \item  \citet{clarke_longevity_2025}
    \item  \citet{calegario_documentation_2021}
    \item  \citet{colazo_licence_2009}
    \item  \citet{perkel_challenge_2020}
    \item  \citet{robert_reproducibility_2020}
    \item  \citet{mcpherson_nimehub_2016}
    \item  \citet{bonvoisin_source_2017}
    \item  \citet{fiordelmondo_longevity_2024}    
\end{enumerate}

\subsection{State of Open Source at NIME 2025}
An audit of NIME 2025 papers and the alignment with FAIR principles.

\section{An Open Methodology for DMIs}
Presentation of an open methodology that can be applied to DMIs using \cite{hamilton_augmentation_2025_activities} as a case study covering the following elements.

\subsection{Source}
Modular source, with though given to dependencies and how parts of the project may be used \textit{outside} of the project itself.
\subsection{Version Control}
A repository to navigate through the evolution of a project.
If a meta repository approach is taken, the meta repository will always be a snapshot, less granular than the individual submodules.
\subsection{Archiving}
The archiving of the project, collecting all elements and assigning a DOI and a space that has longevity. Institutional Archive or independent, like  Zenodo.
\subsection{Version Control $\neq$ Archiving}
A version control repository is not a substitute for an archival repository.

on software and citation

\begin{enumerate}
    \item  \citet{stall_journal_2023}
    \item  \citet{soito_citations_2016}
    \item  \citet{katz_recognizing_2021}
    \item  \citet{smith_software_2016}
\end{enumerate}

\subsection{Documentation}
Diataxis can be used as a framework. The documentation addresses two groups, users and developers. 

\subsection{Licensing}
A suitable open source licence, the choice of which has implication on archiving on platforms like zenodo

\section{Concluding remarks and Recommendations}

If you want to engage with wider communities, reduce the friction for that community to engage with the work.
Readers should not be burdened with detective work and solving the problem of creating a standard disseminating and project source should not be replicate across authors.
Suggested solutions: a field in the metadata and a space in the paper template for 1. the project VCS repository and 2. the project archive.
% A recent report shows a clear digital skill gap in humanities compared to the sciences \cite{taylor_shaping_2022}.